\exercisesection{\S~2}

\begin{enumerate}
\def\labelenumi{\arabic{enumi})}
\tightlist
\item
  Let \(\mathfrak{g}\) be a simple Lie algebra of dimension 3, and \(\varPhi\) its Killing form (cf.~§1, Exerc. 15, 16, 17).
\end{enumerate}

\begin{enumerate}
\def\labelenumi{\alph{enumi})}
\item
  An element \(x \in g\) is regular if and only if \(\Phi(x, x) \neq 0\) . Let \(h_x = kx\) be the Cartan subalgebra generated by such an element. Show that \(h_x\) is splitting if and only if \(2\Phi(x, x)\) is a square in k.
\item
  Show that
\end{enumerate}

\(\mathfrak{g}\) is splittable \(\Longleftrightarrow \mathfrak{g}\) is isomorphic to \(\mathfrak{sl}(2,k)\) .

\begin{enumerate}
\def\labelenumi{\arabic{enumi})}
\setcounter{enumi}{1}
\tightlist
\item
  Let \(k_{1}\) be an extension of \(k\) of finite degree \(n \geq 2\) .
\end{enumerate}

\begin{enumerate}
\def\labelenumi{\alph{enumi})}
\item
  Show that the semi-simple k-algebra \(\mathfrak{sl}(2,k_{1})\) is not splittable.
\item
  Show that the splittable simple k-algebra \(\mathfrak{sl}(n,k)\) contains a Cartan subalgebra \(h_{1}\) that is not splitting. (Choose an embedding of the algebra \(k_{1}\) into \(\mathbf{M}_{n}(k)\) , and take \(\mathfrak{h}_{1}=k_{1}\cap\mathfrak{sl}(n,k)\) .)
\end{enumerate}

\begin{enumerate}
\def\labelenumi{\arabic{enumi})}
\setcounter{enumi}{2}
\tightlist
\item
  Let \((\mathfrak{g},\mathfrak{h})\) be a split semi-simple Lie algebra, R its root system, and K the restriction of the Killing form of \(\mathfrak{g}\) to \(\mathfrak{h}\) . With the notations of Chap. VI, §1, no. 12, we have \(K = B_{R^{\vee}}\) and \(K = 4\gamma(R)\varPhi_{R^{\vee}}\) if R is irreducible; moreover, if all the roots of R have the same length, then
\end{enumerate}

\[
\mathrm{K} (H _ {\alpha}, H _ {\alpha}) = 4 h \quad \text { for   all } \alpha \in \mathrm{R},
\]

where \(h\) is the Coxeter number of \(\mathrm{W}(\mathbb{R})\) .

If for example \(\mathfrak{g}\) is simple of type \(\mathrm{E}_6, \mathrm{E}_7\) or \(\mathrm{E}_8\) , \(\mathrm{K}(H_{\alpha}, H_{\alpha})\) is equal to 48, 72 or 120.

\begin{enumerate}
\def\labelenumi{\arabic{enumi})}
\setcounter{enumi}{3}
\tightlist
\item
  Let \((\mathfrak{g},\mathfrak{h})\) be a split semi-simple Lie algebra, and \((X_{\alpha})_{\alpha \in \mathrm{R}}\) a family of elements satisfying the conditions of Lemma 2. If \(\alpha ,\beta \in \mathrm{R}\) and \(\alpha +\beta \in \mathrm{R}\) , define \(\mathrm{N}_{\alpha \beta}\) by the formula \([X_{\alpha},X_{\beta}] = \mathrm{N}_{\alpha \beta}X_{\alpha +\beta}\) ; if \(\alpha +\beta \notin \mathrm{R}\) , put \(\mathrm{N}_{\alpha \beta} = 0\) . Prove the following formulas:
\end{enumerate}

\(a)\mathrm{N}_{\alpha \beta} = -\mathrm{N}_{\beta \alpha}.\)

\begin{enumerate}
\def\labelenumi{\alph{enumi})}
\setcounter{enumi}{1}
\tightlist
\item
  If \(\alpha, \beta, \gamma \in \mathbb{R}\) are such that \(\alpha + \beta + \gamma = 0\) , then
\end{enumerate}

\[
\frac {\mathrm{N} _ {\alpha \beta}}{\langle \gamma , \gamma \rangle} = \frac {\mathrm{N} _ {\beta \gamma}}{\langle \alpha , \alpha \rangle} = \frac {\mathrm{N} _ {\gamma \alpha}}{\langle \beta , \beta \rangle}.
\]

\begin{enumerate}
\def\labelenumi{\alph{enumi})}
\setcounter{enumi}{2}
\tightlist
\item
  If \(\alpha, \beta, \gamma, \delta \in \mathbb{R}\) are such that \(\alpha + \beta + \gamma + \delta = 0\) , and if no pair has sum zero, we have:
\end{enumerate}

\[
\frac {\mathrm{N} _ {\alpha \beta} \mathrm{N} _ {\gamma \delta}}{\langle \gamma + \delta , \gamma + \delta \rangle} + \frac {\mathrm{N} _ {\beta \gamma} \mathrm{N} _ {\alpha \delta}}{\langle \alpha + \delta , \alpha + \delta \rangle} + \frac {\mathrm{N} _ {\gamma \alpha} \mathrm{N} _ {\beta \delta}}{\langle \beta + \delta , \beta + \delta \rangle} = 0.
\]

\begin{enumerate}
\def\labelenumi{\arabic{enumi})}
\setcounter{enumi}{4}
\item
  Let \((\mathfrak{g},\mathfrak{h})\) be a split semi-simple Lie algebra, and \((X_{\alpha})_{\alpha\in\mathbb{R}}\) a family of elements satisfying the conditions of Lemma 2. Let B be a basis of R. The \(H_{\alpha}\) ( \(\alpha\in B\) ) and the \(X_{\alpha}\) ( \(\alpha\in R\) ) form a basis of the vector space g. Show that the discriminant of the Killing form of g with respect to this basis (Algebra, Chap. IX, §2) is a rational number, independent of the choice of h, of B, and of the \(X_{\alpha}\) .¹⁰ Deduce that, if \(n=\dim g\) , the element of \(\bigwedge^{n}g\) defined by the exterior product of the \(H_{\alpha}\) ( \(\alpha\in B\) ) and the \(X_{\alpha}\) ( \(\alpha\in R\) ) is independent, up to sign, of the choice of h, of B, and of the \(X_{\alpha}\) .

\footnotetext[10]{For an explicit calculation of this discriminant, see: T. A. SPRINGER and R. STEINBERG, Conjugacy Classes (no. 4.8), Seminar on Algebraic Groups and Related Finite Groups, Lect. Notes in Math. 131, Springer-Verlag (1970).}
  \item
  Let \((X_{\alpha})_{\alpha\in\mathbb{R}}\) be a Chevalley system in the split semi-simple Lie algebra \((\mathfrak{g},\mathfrak{h})\) . Let \(\alpha,\beta\in R\) , and let p (resp. q) be the largest integer j such that \(\beta+j\alpha\in R\) (resp. \(\beta-j\alpha\in R\) ), cf.~Lemma 4. Show that
\end{enumerate}

\[
\operatorname{ad} (X _ {\alpha}) ^ {k} (X _ {\beta - q \alpha}) = \pm k! X _ {\beta + (k - q) \alpha} \text { for } 0 \leq k \leq p + q.
\]

Deduce that the algebra \(\mathfrak{g}_{\mathbf{Z}}\) of Prop. 8 is stable under the \(\mathrm{ad}(X_{\alpha})^k / k!\) and under the \(e^{\mathrm{ad}(X_{\alpha})}\) (cf.~§12).

\begin{enumerate}
\def\labelenumi{\arabic{enumi})}
\setcounter{enumi}{6}
\tightlist
\item
  Assume that \(k\) is an ordered field. Let \((\mathfrak{g},\mathfrak{h})\) be a split semi-simple Lie algebra of rank \(l\) ; put \(\dim \mathfrak{g} = l + 2m\) .
\end{enumerate}

\begin{enumerate}
\def\labelenumi{\alph{enumi})}
\item
  Show that the Killing form \(\varPhi\) of g is the direct sum of a neutral form of rank 2m and a positive non-degenerate form of rank l; in particular, its index is m.
\item
  Let \(\varphi\) be an involutive automorphism of \(\mathfrak{g}\) whose restriction to \(\mathfrak{h}\) is -Id. Show that the form
\end{enumerate}

\[
(x, y) \mapsto - \Phi (x, \varphi (y)), \quad x, y \in \mathfrak {g},
\]

is symmetric, non-degenerate, and positive.

\(c)\) Let \(k' = k(\sqrt{\alpha})\) , where \(\alpha\) is an element \(< 0\) in \(k\) . Denote by \(c\) the non-trivial \(k\) -automorphism of \(k'\) . Let \(\mathfrak{g}_c\) be the \(k\) -subspace of \(\mathfrak{g}_{(k')} = k' \otimes_k \mathfrak{g}\) formed by the elements \(y\) such that

\[
(1 \otimes \varphi). y = (c \otimes 1). y.
\]

Show that \(g_{c}\) is a k-Lie subalgebra of \(\mathfrak{g}_{(k^{\prime})}\) and that the injection of \(g_{c}\) into \(\mathfrak{g}_{(k^{\prime})}\) extends to an isomorphism from \(k^{\prime}\otimes_{k}g_{c}\) to \(\mathfrak{g}_{(k^{\prime})}\) . The algebra \(g_{c}\) is semi-simple, and \(\sqrt{\alpha}h\) is a Cartan subalgebra of it.

\(d\) ) Show that the Killing form of \(\mathfrak{g}_c\) is negative. Deduce that \(\mathfrak{g}_c\) is not splittable (unless \(\mathfrak{g} = 0\) ).

\begin{enumerate}
\def\labelenumi{\alph{enumi})}
\setcounter{enumi}{4}
\tightlist
\item
  When \(k = \mathbf{R}\) , show that \(\operatorname{Int}(\mathfrak{g}_c)\) is compact.
\end{enumerate}

\begin{enumerate}
\def\labelenumi{\arabic{enumi})}
\setcounter{enumi}{7}
\tightlist
\item
  \begin{enumerate}
  \def\labelenumii{\alph{enumii})}
  \tightlist
  \item
    Let g be a Lie algebra and n an integer \(\geq 0\) . Let \(\Sigma_{n}(\mathfrak{g})\) be the subset of \(g^{n}\) consisting of the families \((x_{1},\ldots,x_{n})\) that generate g as a k-algebra. Show that \(\Sigma_{n}(\mathfrak{g})\) is open in \(g^{n}\) in the Zariski topology (Chap. VII, App. I). If \(k'\) is an extension of k, then \(\Sigma_{n}(\mathfrak{g}_{(k')})\cap\mathfrak{g}^{n}=\Sigma_{n}(\mathfrak{g})\) . Deduce that, if \(\mathfrak{g}_{(k')}\) can be generated by n elements, so can g.
  \end{enumerate}
\end{enumerate}

\begin{enumerate}
\def\labelenumi{\alph{enumi})}
\setcounter{enumi}{1}
\item
  Let \((\mathfrak{g},\mathfrak{h})\) be a split semi-simple Lie algebra. Let \(x\) be an element of \(\mathfrak{h}\) such that \(\alpha (x)\neq 0\) for all \(\alpha \in \mathrm{R}\) and \(\alpha (x)\neq \beta (x)\) for any pair of distinct elements \(\alpha ,\beta \in \mathrm{R}\) . For all \(\alpha \in \mathrm{R}\) , let \(y_{\alpha}\) be a non-zero element of \(\mathfrak{g}^{\alpha}\) , and let \(y = \sum_{\alpha \in \mathrm{R}}y_{\alpha}\) . Show that, for all \(\alpha \in \mathrm{R}\) , there exists a polynomial \(\mathrm{P}_{\alpha}(\mathrm{T})\in k[\mathrm{T}]\) , without constant term, such that \(y_{\alpha} = \mathrm{P}_{\alpha}(\mathrm{ad}x).y\) . Deduce that \(\mathfrak{g}\) is generated by \(\{x,y\}\) .
\item
  Show, by means of \(a\) ) and \(b\) ), that every semi-simple Lie algebra can be generated by two elements.
\end{enumerate}

¶9) Let G be a connected finite dimensional real Lie group. Let g be its Lie algebra, and let \(\{x_{1},\ldots,x_{n}\}\) be a generating family of g. For \(m\geq0\) , denote by \(\Gamma_{m}\) the subgroup of G generated by the \(\exp(2^{-m}x_{i})\) , \(1\leq i\leq n\) . We have \(\Gamma_{0}\subset\Gamma_{1}\subset\cdots\) .

\begin{enumerate}
\def\labelenumi{\alph{enumi})}
\item
  Show that the union of the \(\Gamma_{m}\) is dense in G.
\item
  Let \(H_{m}\) be the identity component of the closure \(\overline{\Gamma}_{m}\) of \(\Gamma_{m}\) . We have \(H_{0} \subset H_{1} \subset \cdots\) , and the family ( \(H_{m}\) ) is stationary; let H be the common value of the \(H_{m}\) for m sufficiently large. Show that H is normal in G (observe that H is normalized by all the \(\Gamma_{m}\) ), and that the image of \(\Gamma_{m}\) in G/H is a discrete subgroup of G/H. The union of these subgroups being dense in G/H, deduce (cf.~Chap. III, §6, Exerc. 23 d)) that G/H is nilpotent.
\item
  Assume that \(\mathfrak{g} = \mathcal{D}\mathfrak{g}\) . Show that \(\mathrm{G} = \mathrm{H}\) , in other words, that \(\Gamma_{m}\) is dense in \(\mathrm{G}\) if \(m\) is sufficiently large.
\end{enumerate}

\begin{enumerate}
\def\labelenumi{\arabic{enumi})}
\setcounter{enumi}{9}
\item
  Let \(G\) be a connected semi-simple real Lie group. Show, by using Exerc. 8 and 9, that there exists a dense subgroup of \(G\) generated by two elements.
\item
  Let \(\mathrm{R}\) be a root system of rank \(l\) , and \(\varPhi_{\mathrm{R}}\) the corresponding canonical bilinear form (Chap. VI, §1, no. 12). The matrix \(\varPhi = (\varPhi_{\mathrm{R}}(\alpha, \beta))_{\alpha, \beta \in \mathrm{R}}\) is of rank \(l\) , and \(\varPhi^2 = \varPhi\) . Deduce the formula \(\sum_{\alpha \in \mathrm{R}} \varPhi_{\mathrm{R}}(\alpha, \alpha) = \operatorname{Tr} \varPhi = l\) .
\end{enumerate}

